\begin{center}
\textsc{Gauss} to \textsc{Olbers}.\quad Göttingen, 2~April 1826.
\end{center}

\par\noindent\hspace*{\parindent}\mbox{.\ .\ .\ .\ .\ .}\quad As for my survey, there is little I can decide \textls[120]{on my own}.
My \textls[120]{commission} is essentially complete as far as the trigonometric part is concerned, and I thought that, if \textsc{Schumacher} is willing to give me \textls[120]{adequate} support, I would make the zenith-sector observations in late summer.
The precarious state of my health deters me from proposing more extensive trigonometric measurements. In the meantime, \textls[120]{between ourselves}, I recently wrote to \textsc{Münster} to say that I was willing to devote my efforts to them in the future as well, should such measurements be requested.

\par\noindent\hspace*{\parindent}My theoretical work, with its \textls[120]{very} wide scope, unfortunately still has many gaps. The easiest way out for me would be if I
\par\clearpage
\noindent\null\hfill{\small\textbf{ON THE HANOVER TRIANGULATION.}}\hfill\makebox[0pt][r]{\fontsize{13}{16}\selectfont \textbf{377}}\par\medskip
\noindent allowed myself to include all my computational procedures when publishing my measurements, but reserved their \textls[120]{derivation} from more fundamental principles for an entirely separate work in happier times to come.
Then there would be no difficulty anywhere.
For the present I shall undertake the precise adjustment of my 32 points, which give 51 triangles and 146 directions.
I have completed the adjustment of the heights (a very much easier task) in the past few days.\quad\mbox{.\ .\ .\ .\ .\ .}
